\documentclass[10pt]{article}
\usepackage[T1]{fontenc}
\usepackage[utf8]{inputenc}
\usepackage[french]{babel}
\usepackage{amsmath,amssymb,mathrsfs,mathpazo}
\usepackage{geometry}
\geometry{a4paper,margin=2cm}
\usepackage{graphicx,array,tabularx,tabularray,colortbl,multirow,multicol,booktabs,enumitem,pifont,ulem,fancybox,diagbox,cancel}
\usepackage{pstricks,pstricks-add,pst-plot,pst-tree,pst-node,pst-text,pst-func,pst-math,pst-bspline,pst-3dplot,pst-eucl,pst-solides3d}
\let\pscalculate\undefined
\usepackage{pst-bezier}
\usepackage{tikz,pgfplots,tkz-tab}
\usetikzlibrary{arrows,arrows.meta,patterns,calc,positioning,trees,decorations.markings,shapes.geometric,shapes.misc,shapes.arrows,intersections,decorations.pathmorphing,decorations.pathreplacing,backgrounds,shadows,babel}
\pgfplotsset{compat=1.18}
\usepgfplotslibrary{fillbetween}
\usepackage[np]{numprint}
\usepackage{listings,makecell,pas-tableur,twoopt}
\let\barre\undefined
\usepackage[pstricks]{tablvar}
\definecolor{darkgreen}{RGB}{0,100,0}
\providecommand{\red}{\color{red}}
\providecommand{\blue}{\color{blue}}
\providecommand{\green}{\color{green}}
\providecommand{\black}{\color{black}}
\providecommand{\white}{\color{white}}
\providecommand{\gray}{\color{gray}}
\providecommand{\cyan}{\color{cyan}}
\providecommand{\magenta}{\color{magenta}}
\providecommand{\yellow}{\color{yellow}}
\providecommand{\violet}{\color{violet}}
\providecommand{\orange}{\color{orange}}
\providecommand{\euro}{\texteuro}
\providecommand{\EUR}{\texteuro}
\providecommand{\ieme}{\textsuperscript{e}}
\providecommand{\degres}{\ensuremath{^\circ}}
\providecommand{\degree}{\ensuremath{^\circ}}
\providecommand{\vv}[1]{\overrightarrow{#1}}
\providecommand{\Coord}[1]{#1}
\providecommand{\decosix}{}
\providecommand{\starredbullet}{\ensuremath{\bullet}}
\pagestyle{empty}
\setlength{\parindent}{0pt}
\usepackage{adjustbox}\DeclareUnicodeCharacter{2236}{:}\DeclareUnicodeCharacter{2190}{\ensuremath{\leftarrow}}
\providecommand{\R}{\mathbb{R}}
\providecommand{\N}{\mathbb{N}}
\providecommand{\Z}{\mathbb{Z}}
\providecommand{\Q}{\mathbb{Q}}
\providecommand{\D}{\mathbb{D}}
\providecommand{\C}{\mathbb{C}}
\providecommand{\e}{{\mathrm{e}}}
\providecommand{\vect}[1]{\overrightarrow{#1}}
\providecommand{\vectt}[1]{\overrightarrow{\mathrm{#1}}}
\providecommand{\vvt}[1]{\overrightarrow{\mathrm{#1}}}
\providecommand{\np}[2][]{{#2~#1}}
\providecommand{\vv}[1]{{\overrightarrow{#1}}}
\renewcommand{\d}{{\mathrm{d}}}
\providecommand{\degres}{{^\circ}}
\providecommand{\degre}{{^\circ}}
\providecommand{\blue}{\color{blue}}
\providecommand{\red}{\color{red}}
\providecommand{\enskip}{\quad}
\providecommand{\llbracket}{\lbrack\!\lbrack}
\providecommand{\rrbracket}{\rbrack\!\rbrack}
\providecommand{\euro}{{\text{€}}}
\providecommand{\no}{{\text{n°}}}
\providecommand{\interval}[2]{{[#1;#2]}}
\providecommand{\intervFO}[2]{{[#1;#2[}}
\providecommand{\dotfill}{\ldots}

\providecommand{\ds}{\displaystyle}
\providecommand{\textding}[1]{#1}
\providecommand{\Oij}{\(\mathrm{O};\vec{\imath},\vec{\jmath}\)}
\providecommand{\Oijk}{\(\mathrm{O};\vec{\imath},\vec{\jmath},\vec{k}\)}
\providecommand{\Ouv}{\(\mathrm{O};\vec{u},\vec{v}\)}
\providecommand{\intervFF}[2]{\ensuremath{[#1;#2]}}
\renewcommand{\euro}{\ifmmode\text{\texteuro}\else\texteuro\fi}
\DeclareUnicodeCharacter{00B0}{\ifmmode^{\circ}\else\textdegree\fi}
\DeclareUnicodeCharacter{20AC}{\ifmmode\text{\texteuro}\else\texteuro\fi}
% Énoncés : Éducation nationale. Transcription et corrigés : APMEP, Denis Vergès et contributeurs. Découpage et classement : CoopMaths, Sébastien Lozano, Éric Elter et contributeurs. Adaptation de la présentation : Maths Explorer.
% CC BY-SA : https://coopmaths.fr/www/about/
% https://coopmaths.fr/alea/static/eam/2026/tex/eam_technologique_2026_06_antilles_automatismes.tex
\begin{document}
\noindent\textit{APMEP / CoopMaths. CC BY-SA. Présentation adaptée par Maths Explorer.}\par\medskip
\textbf{Pour cette première partie, aucune justification n'est demandée et une seule réponse est possible par question.\\
Pour chaque question, reportez son numéro sur votre copie et indiquez votre réponse.}

Une réponse fausse ou l'absence de réponse n'enlève aucun point.

\subsection*{Question 1}
	Un objet coûtait 100 \euro{}.

	Après une remise, il coûte 80 \euro{}.

	Quel était le pourcentage de cette remise ?

	
	\medskip
	\begin{tblr}{colspec={@{}*{4}{X[c]}}}
		\textbf{A.} $0,2\,\%$ &
		\textbf{B.} $20\,\%$  &
		\textbf{C.} $0,8\,\%$ &
		\textbf{D.} $80\,\%$
	\end{tblr}

\subsection*{Question 2}
	Dans une salle de cinéma, $60\,\%$ des personnes ont moins de 25 ans, dont $10\,\%$ sont mineurs.

	Quelle est la proportion de mineurs dans la salle ?

	\medskip
	\begin{tblr}{colspec={@{}*{4}{X[c]}}}
		\textbf{A.} 0,1  &
		\textbf{B.} 0,5  &
		\textbf{C.} 0,06 &
		\textbf{D.} 0,6
	\end{tblr}

\subsection*{Question 3}
	Un prix augmente de $10\,\%$ puis diminue de $20\,\%$.

	Quelle est l'évolution globale de ce prix ?

	\medskip
	\begin{tblr}{colspec={@{}*{4}{X[c]}}}
		\textbf{A.} Baisse de 10\,\% &
		\textbf{B.} Baisse de 12\,\% &
		\textbf{C.} Hausse de 10\,\% &
		\textbf{D.} Baisse de 88\,\%
	\end{tblr}

\bigskip

\textbf{Pour les questions 4 et 5}, on considère la formule suivante permettant de transformer des degrés Celsius (\textdegree{}C) en degrés Fahrenheit (\textdegree{}F) :
\[F=1,8 C+32\]
où $F$ désigne la température en \textdegree{}F et $C$ désigne la température en \textdegree{}C.

\subsection*{Question 4}
	Sachant que l'eau bout à 100~\textdegree{}C, la température d'ébullition de l'eau en \textdegree{}F est :

	\medskip
	\begin{tblr}{colspec={@{}*{4}{X[c]}}}
		\textbf{A.} 33,8~\textdegree{}F &
		\textbf{B.} 50~\textdegree{}F   &
		\textbf{C.} 100~\textdegree{}F  &
		\textbf{D.} 212~\textdegree{}F
	\end{tblr}

\subsection*{Question 5}
	La formule permettant de transformer des degrés Fahrenheit (\textdegree{}F) en degrés Celsius (\textdegree{}C) est donc :

	\medskip
	\begin{tblr}{colspec={@{}*{4}{X[c]}}}
		\textbf{A.} $C=\dfrac{F}{1,8}-32$ &
		\textbf{B.} $C=\dfrac{F-32}{1,8}$ &
		\textbf{C.} $C=\dfrac{F-1,8}{32}$ &
		\textbf{D.} $C=\dfrac{F}{32}-1,8$ \\
	\end{tblr}

\subsection*{Question 6}
	On considère $A=\dfrac{~~\dfrac{5}{3}~~}{15}$ :

	\medskip
	\begin{tblr}{colspec={@{}*{4}{X[c]}}}
		\textbf{A.} $A = \dfrac{1}{9}$  &
		\textbf{B.} $A = 25$            &
		\textbf{C.} $A = \dfrac{1}{45}$ &
		\textbf{D.} $A = 5$
	\end{tblr}

\subsection*{Question 7}
	Dans le repère ci-dessous, on a représenté une droite $d$.

	\begin{center}
		\begin{tikzpicture}[x=5mm,y=5mm,> = stealth]
			\draw[gray,line cap = round,line width=0.25mm,dash pattern= on 0mm off 1mm,xstep=1,ystep=1] (-3,-1.2) grid (9.5,4.8);
			\draw[thick,->] (-3,0)--(9.8,0) node[below left]{$x$};
			\draw[thick,->] (0,-1.3)--(0,4.8) node[below left]{$y$};
			\foreach \x in {-3,...,9}
			{\draw[thick] (\x,2pt)--(\x,-2pt);}
			\foreach \y in {-1,1,2,3,4}
			{\draw[thick] (2pt,\y)--(-2 pt,\y);}
			\node[below left] at (0,0) {0};
			\node[below ] at (1,0) {1};
			\node[left ] at  (0,1) {1};
			\clip (-3,-1.3) rectangle (8.5,4.8);
			\draw[red,thick] (-4,4)--(8,-2);
			\node[red] at(-1.5,3.5) {$(d)$};
		\end{tikzpicture}
	\end{center}

	L'équation réduite de cette droite est :

	\medskip
	\begin{tblr}{colspec={@{}*{4}{X[c]}}}
		\textbf{A.} $y=-0,5 x+4$ &
		\textbf{B.} $y=-0,5 x+2$ &
		\textbf{C.} $y=-2 x+4$   &
		\textbf{D.} $y=-2 x+2$
	\end{tblr}

\subsection*{Question 8}
	Dans un repère, on considère les points E de coordonnées $(20 ; 25)$ et F de coordonnées $(5 ; 15)$.

	Le coefficient directeur de la droite (EF) est alors égal à :

	\medskip
	\begin{tblr}{colspec={@{}*{4}{X[c]}}}
		\textbf{A.} 2              &
		\textbf{B.} $\dfrac{3}{2}$ &
		\textbf{C.} $\dfrac{2}{3}$ &
		\textbf{D.} $\dfrac{1}{2}$
	\end{tblr}

\subsection*{Question 9}
	Dans le repère ci-dessous, on a représenté une fonction $f$ définie sur l'intervalle $[-5 ; 3]$.

	\begin{center}
		\begin{tikzpicture}[x=5mm,y=5mm,> = stealth]
			\draw[gray,line cap = round,line width=0.25mm,dash pattern= on 0mm off 1mm,xstep=1,ystep=1] (-5.8,-3.4) grid (4.6,3.8);
			\draw[thick,-latex] (-5.8,0)--(4.6,0) node[below left]{$x$};
			\draw[thick,-latex] (0,-3.4)--(0,3.8) node[below left]{$y$};
			\foreach \x in {-5,...,3}
			{\draw[thick] (\x,2pt)--(\x,-2pt);}
			\foreach \y in {-3,...,3}
			{\draw[thick] (2pt,\y)--(-2 pt,\y);}
			\node[below left] at (0,0) {0};
			\node[below ] at (1,0) {1};
			\node[left ] at  (0,1) {1};

			\draw[teal,line width=1.2pt,smooth]
			(-5,0.5) .. controls (-4.5,1.5) and (-3.8,2.5) ..
			(-3,2.5) .. controls (-2.2,2.5) and (-1.5,1) ..
			(-1,0) .. controls (-0.5,-1) and (0.5,-2.5) ..
			(1.5,-2.5) .. controls (2,-2.5) and (2.5,-1.75) ..
			(3,-1);
		\end{tikzpicture}
	\end{center}

	Le tableau de signes de cette fonction est donné par :

	\begin{tblr}{colspec={@{}*{2}{X}}}
		\textbf{A.}~ \begin{tikzpicture}[baseline={(F0.base)}]
			\tkzTabInit[lgt=2.2,espcl=1.2,deltacl=0.4]
			{$x$/0.5, signe de $f(x)$/1}{$-5$,$-3$,$\np{1.5}$,$3$}
			\tkzTabLine{,+,z,-,z,+,}
		\end{tikzpicture} &
		\textbf{B.}~ \begin{tikzpicture}[baseline={(F0.base)}]
			\tkzTabInit[lgt=2.2,espcl=1.8,deltacl=0.4]
			{$x$/0.5, signe de $f(x)$/1}{$-5$,$-1$,$3$}
			\tkzTabLine{,+,z,-,}
		\end{tikzpicture}\\
			\textbf{C.}~ \begin{tikzpicture}[baseline={(F0.base)}]
				\tkzTabInit[lgt=2.2,espcl=1.8,deltacl=0.4]
				{$x$/0.5, signe de $f(x)$/1}
				{$-5$,$\np{-1.5}$,$3$}
				\tkzTabLine{,-,z,+,}
			\end{tikzpicture} &
		\textbf{D.}~ \begin{tikzpicture}[baseline={(F0.base)}]
			\tkzTabInit[lgt=2.2,espcl=1.2,deltacl=0.4]
			{$x$/0.5, signe de $f(x)$/1}{$-5$,$-1$,$0$,$3$}
			\tkzTabLine{,+,z,-,z,+,}
		\end{tikzpicture}
	\end{tblr}


\subsection*{Question 10}
	Les paramètres statistiques d'une série de notes d'un contrôle noté sur 20 sont donnés dans le tableau ci-dessous.
	\begin{center}
		\begin{tblr}{colspec={X[c]X[c]X[c]X[1.15,c]X[c]},vlines}
			\hline
			Minimum & Premier quartile & Médiane & Troisième quartile & Maximum \\ \hline
			2       & 5                & 10      & 12                 & 17      \\ \hline
		\end{tblr}
	\end{center}

	La proportion d'élèves ayant une note inférieure ou égale à 12 est :

	\begin{tblr}{colspec={@{}*{2}{X}}}
		\textbf{A.} inférieure ou égale à 25\,\% &
		\textbf{B.} égale à 50\,\% \\
		\textbf{C.} supérieure ou égale à 75\,\% &
		\textbf{D.} égale à 100\,\%
	\end{tblr}

\subsection*{Question 11}
	Le diagramme ci-dessous donne la répartition en pourcentage de joueurs de jeux vidéo selon leur âge.

	\begin{center}
	\begin{tikzpicture}
		\draw[black,fill=gray1] (0,0)--(90:2.25) arc (90:59.4:2.25) ;
		\draw[black,fill=gray9] (0,0)--(59.4:2.25) arc (59.4:14.4:2.25) ;
		\draw[black,fill=gray6] (0,0)--(14.4:2.25) arc (14.4:-53.64:2.25) ;
		\draw[black,fill=gray3] (0,0)--(-53.64:2.25) arc (-53.64:-120.6:2.25) ;
		\draw[black,fill=gray5] (0,0)--(-120.6:2.25) arc (-120.6:-209.88:2.25) ;
		\draw[black,fill=gray8] (0,0)--(-209.88:2.25) arc (-209.88:-270:2.25)--cycle ;
		\draw (74.7:2.3)--(74.7:2.8)-- ++(0.2,0) node[right,text width=20mm,align = center]{6--9~ans 8,5};
		\draw (36.9:2.3)--(36.9:2.8)-- ++(0.2,0) node[right,text width=20mm,align = center]{10--14~ans 12,5};
		\draw (-19.6:2.3)--(-19.6:2.8)-- ++(0.2,0) node[right,text width=20mm,align = center]{15--24~ans 18,9};
		\draw (-87.1:2.3)--(-87.2:2.8)-- ++(0.2,0) node[right,text width=20mm,align = center]{25--34~ans 18,6};
		\draw (-165.2:2.3)--(-165.2:2.8)-- ++(-0.2,0) node[left,text width=20mm,align = center]{35--49~ans 24,8};
		\draw (-240:2.3)--(-240:2.8)-- ++(-0.2,0) node[left,text width=20mm,align = center]{50--65~ans 16,7};
	\end{tikzpicture}

	\emph{Source : CNC -- TNS Sofres 2014}
	\end{center}

	La proportion de joueurs de 24 ans et moins est :

	\medskip
	\begin{tblr}{colspec={@{}*{4}{X[c]}}}
		\textbf{A.} 60,1\,\%   &
		\textbf{B.} $18,9\,\%$ &
		\textbf{C.} $39,9\,\%$ &
		\textbf{D.} $\dfrac{1}{3}$
	\end{tblr}

\subsection*{Question 12}
	Une enquête réalisée auprès des 800 élèves d'un lycée donne les résultats suivants :

	\begin{center}
		\begin{tblr}{colspec={X[3,c]X[c]X[c]X[c]},hlines, vlines}
			L'élève                                 & est une fille & est un garçon & Total \\
			ne pratique aucune activité sportive    & 90            & 200           & 290   \\
			pratique au moins une activité sportive & 210           & 300           & 510   \\
			Total                                   & 300           & 500           & 800   \\
		\end{tblr}
	\end{center}

	On choisit au hasard un élève parmi les garçons.

	Quelle est la probabilité qu'il ne pratique aucune activité sportive ?

	\medskip
	\begin{tblr}{colspec={@{}*{4}{X[c]}}}
		\textbf{A.} $\dfrac{200}{290}$&
		\textbf{B.} $\dfrac{200}{800}$&
		\textbf{C.} $\dfrac{290}{800}$&
		\textbf{D.} $\dfrac{200}{500}$
	\end{tblr}

\end{document}
